\oldpage[0\textsubscript{I}]{35}

\subsection{Haces de espacios pseudodiscretos}

\begin{env}[3.8.1]
\label{0.3.8.1-fr}
Sea $X$ un espacio topológico cuya topología admite una base
$\mathfrak{B}$ formada por conjuntos abiertos \emph{cuasicompactos}. Sea
$\mathcal{F}$ un \emph{haz de conjuntos} sobre $X$. Si se dota a cada
$\mathcal{F}(U)$ de la topología \emph{discreta},
$U\to\mathcal{F}(U)$ es un \emph{prehaz de espacios topológicos}.
Vamos a ver que existe un \emph{haz de espacios topológicos
$\mathcal{F}'$ asociado} a $\mathcal{F}$ (3.5.6) tal que
$\Gamma(U,\mathcal{F}')$ es el espacio discreto $\mathcal{F}(U)$ para todo
abierto \emph{cuasicompacto} $U$. Para ello bastará mostrar que el
prehaz $U\to\mathcal{F}(U)$ de espacios topológicos discretos
\emph{sobre $\mathfrak{B}$} verifica la condición $(\mathrm{F}_0)$ de (3.2.2)
y, más generalmente, que si $U$ es un abierto cuasicompacto y
$(U_\alpha)$ es un recubrimiento de $U$ por conjuntos de $\mathfrak{B}$,
la topología menos fina $\mathcal{T}$ sobre $\Gamma(U,\mathcal{F})$ que hace
continuas las aplicaciones
$\Gamma(U,\mathcal{F})\to\Gamma(U_\alpha,\mathcal{F})$ es la topología
\emph{discreta}. Ahora bien, hay un número finito de índices $\alpha_i$ tales
que $U=\bigcup_i U_{\alpha_i}$. Sea $s\in\Gamma(U,\mathcal{F})$ y sea $s_i$
su imagen en $\Gamma(U_{\alpha_i},\mathcal{F})$. La intersección de las
imágenes inversas de los conjuntos $\{s_i\}$ es, por definición, una vecindad
de $s$ para $\mathcal{T}$. Como $\mathcal{F}$ es un haz de conjuntos y los
$U_{\alpha_i}$ recubren $U$, esta intersección se reduce a $s$, lo cual
demuestra nuestra afirmación.

Se observará que, si $U$ es un abierto no cuasicompacto de $X$, el espacio
topológico $\Gamma(U,\mathcal{F}')$ sigue teniendo a $\Gamma(U,\mathcal{F})$
como conjunto subyacente. Su topología, que en general no es discreta, es
la menos fina que hace continuas las aplicaciones
$\Gamma(U,\mathcal{F})\to\Gamma(V,\mathcal{F})$, para $V\in\mathfrak{B}$ y
$V\subset U$ (siendo discretos los $\Gamma(V,\mathcal{F})$).

Las consideraciones precedentes se aplican sin modificación a los haces de
grupos o de anillos (no necesariamente conmutativos), y les asocian
respectivamente haces de \emph{grupos topológicos} o de \emph{anillos
topológicos}. Para abreviar, diremos que el haz $\mathcal{F}'$ es el haz de
\emph{espacios} (resp. de \emph{grupos}, de \emph{anillos})
\emph{pseudodiscretos} asociado al haz de conjuntos (resp. de grupos,
de anillos) $\mathcal{F}$.
\end{env}
